\subsection{Kompilieren einer PDF-Datei}
\label{latex-kompilieren}
Auf dem Arbeitsrechner des Fachschaftsrats werden Dokumente in VSCodium mit der Erweiterung LaTeX Workshop kompiliert, die dort bereits für LuaLaTeX eingerichtet ist:

\begin{enumerate}
	\item Die \texttt{.tex}-Datei in VSCodium öffnen.
	\item Das Dokument mit \texttt{Strg+Alt+B} oder über den grünen Pfeil \textit{Build LaTeX project} oben rechts kompilieren. Standardmäßig wird das Dokument außerdem bei jedem Speichern automatisch neu kompiliert.
	\item Die erzeugte PDF-Datei mit \texttt{Strg+Alt+V} oder über das Symbol \textit{View LaTeX PDF} oben rechts in einer Vorschau neben dem Quelltext öffnen. Die Vorschau aktualisiert sich nach jedem Kompilieren automatisch.
\end{enumerate}
% TODO: Prüfen, ob Tastenkürzel und automatisches Kompilieren beim Speichern auf tesla so eingerichtet sind.

Die PDF-Datei wird im automatisch erstellten Ordner \texttt{build/} abgelegt. Die fertige PDF-Datei kann anschließend beliebig verschoben werden.
